\section{Cycle traces and contraction by the transposition Casimir}\label{sec:casimir}
The estimates in this section keep track of the number of boxes below the
first row of a representation.  The relevant scale is
\[
 h_n(k)=k\log(en/k),\qquad 1\le k\le n.
\]
We will prove a density bound for every injection and a contraction bound
with this same scale, then combine them to control a fixed trace moment.
This gives a complete additional route from the switching network to the
full permutation law.

Let $n=2^d$, and let $\mu_n$ be the law of one directed coordinate sweep.
In any unitary representation its average is
$T_n=\Pi_d\cdots\Pi_1$, where $\Pi_i$ averages the independent fair
switches in direction $i$.  Put $K_n=T_n^*T_n$ and write $\nu_n$ for its
permutation law.  Chronologically, $\nu_n$ runs a sweep and then its
reflection, with fresh independent switches.  For $n=2m$, the network
consists of an outer pair layer, two independent networks with law
$\nu_m$, and another outer pair layer.  We write $K_n(\lambda)$ for the
operator in the irreducible representation $V_\lambda$, and
$D_\lambda=f^\lambda=\dim V_\lambda$.  The same group-algebra operator
also acts on tensor spaces and on injections; its indicated carrier will
always specify which action is meant.

Write $X_{n,k}=[n]_{\ne}^k$ for the ordered distinct $k$-tuples of positions,
and $u_{n,k}$ for their uniform probability measure.  There are
$(n)_k=n(n-1)\cdots(n-k+1)$ such tuples.  For $x,y\in X_{n,k}$ define
\[
 K_{n,k}(x,y)=\Pr_{\pi\sim\nu_n}(\pi x=y),\qquad
 R_x(y)=(n)_kK_{n,k}(x,y),
\]
where $\pi x$ means applying $\pi$ to every entry of $x$.  Thus $R_x$ is
a density relative to $u_{n,k}$.  In particular its $L^p$ norms use a
probability measure, whereas all matrix traces and Schatten norms below
are unnormalized.

\begin{theorem}[Cycle density and Casimir contraction]\label{cas:main}
There are absolute constants $a,C>0$ such that for every dyadic $n$,
every $1\le k\le n$, and every $x\in X_{n,k}$,
\[
 \|R_x\|_{L^{65/64}(u_{n,k})}\le e^{C h_n(k)}.
\]
For every nontrivial partition $\lambda\vdash n$, put $k=n-\lambda_1$.
Then
\[
 \|K_n(\lambda)\|_{\op}\le e^{-a h_n(k)},\qquad
 D_\lambda\Tr K_n(\lambda)^{65}\le e^{C h_n(k)}.
\]
Consequently there is an absolute integer $l$ for which
$\|\mu_n^{*l}-U_{S_n}\|_{\TV}\to0$ as $n\to\infty$ through powers of two.
Here $U_{S_n}$ is the uniform permutation law and total variation includes
the factor $1/2$.  The same bound holds from every deterministic initial
permutation.  When $n=2^d$, these $l$ sweeps take $ld$ physical shuffles.
\end{theorem}

The density proof counts cycles through positive tensor traces.  For the
operator bound we use the transposition Casimir, the sum of the positive
operators $I-U((i\,j))$.  Its scalar on $V_\lambda$ is comparable to $nk$.
Two estimates for a random split of this sum control, respectively, the
loss of total Casimir mass and the imbalance between the two children.
They allow an induction with a fixed multiplicative margin.  Finally,
Hausdorff--Young turns the injection-density bound into the $65$th moment,
and Schatten H\"older applies that moment to powers of the directed sweep.

\subsection{Cycle traces and tuple spreading}
Our first objective is the $L^{65/64}$ density bound.  A selected cycle at
a network node contributes a factor two to the number of possible
routings.  We therefore need an exponential moment for the sum of these
cycle counts.  Positive tensor traces first bound the full cycle count;
a conditional estimate across levels then controls the selected cycles.

If \(\pi\) is a permutation let \(\kappa(\pi)\) denote the number of cycles, counting fixed points. Define for real \(x\ge 1\)
\[
Z_b(x)=\mathbb E_{\pi\sim\nu_b} x^{\kappa(\pi)}
\]
where \(b\) is a power of 2. We first show that there are absolute \(C<\infty\) and \(\gamma>0\) such that, for \(b=2^w\),
\begin{equation}
\log Z_b(1+2^{\lfloor w/8\rfloor})\le C b^{1-\gamma}.
\label{cas:e1}
\end{equation}
For integer \(q\ge 1\), \(Z_b(q)\) is the trace of \(K_b\) on the tensor space \((\mathbb C^q)^{\otimes b}\) with permutations acting by permuting tensor legs. In fact the trace of a permutation counts its fixed words on an alphabet of size \(q\). The following identity is the common tensor calculation underlying this
proof and the harmonic cycle estimates.

\begin{lemma}[Tensor-swap recursion]\label{cas:swap-recursion}
For a dyadic $m$ and an integer $q\ge1$, let $D$ be the action of $K_m$
on $(\mathbb C^q)^{\otimes m}$.  For $S\subseteq[m]$, let
$A_S=\operatorname{Tr}_{S^c}D$, where the partial trace is unnormalized.
Then
\[
 Z_{2m}(q)=2^{-m}\sum_{S\subseteq[m]}\Tr A_S^2,
 \qquad
 Z_{2m}(q)\le Z_m(q)Z_m(1+(q-1)/2).
\]
The scalar function $Z_m(x)$ in the second formula is defined for every
real $x\ge1$, including a nonintegral argument.
\end{lemma}
\begin{proof}
In this representation the recursion at size \(2m\) gives the operator
\[
P(D\otimes D)P,\qquad D=K_m\ \text{on }(\mathbb C^q)^{\otimes m},
\]
where \(P=\prod_{i=1}^m \big((I+F_i)/2\big)\) with \(F_i\) interchanging leg \(i\) between the two copies (number the pair indices by \(i=1,\ldots,m\), using the same ordering in both copies). Here \(P^2=P\) and \(D\) is positive semidefinite. Taking a cyclic trace and expanding \(P\) gives
\begin{equation}
Z_{2m}(q)=2^{-m}\sum_{S\subseteq [m]}\operatorname{Tr}\big[(\operatorname{Tr}_{S^c} D)^2\big],
\label{cas:e2}
\end{equation}
where \(\operatorname{Tr}_{S^c}\) is ordinary, unnormalized partial trace over legs in the complement. Indeed for the product of swaps on \(S\), first trace out the non-swapped legs and then use \(\operatorname{Tr}(F(A\otimes A))=\operatorname{Tr}(A^2)\) with \(F\) the whole swap of the remaining copies.

The operator \(A_S=\operatorname{Tr}_{S^c}D\) is positive semidefinite, with trace \(Z_m(q)\). For an individual permutation operator in \(D\), its partial trace is \(q^{\kappa_0(\pi,S)}\) times a permutation operator on the legs in \(S\), where \(\kappa_0\) counts cycles disjoint from \(S\). This follows by writing matrix entries in the word basis: an input color index is equated to the output index at the next site along the cycle, and input is set equal to output on legs traced out and summed. Thus each deleted whole cycle has a free common index, and on other cycles the deleted legs simply shorten the cycle to its undeleted legs. It follows that \(\|A_S\|\le\mathbb E_{\pi\sim\nu_m} q^{\kappa_0(\pi,S)}\), where \(\|\cdot\|\) on operators denotes operator norm. Use \(\operatorname{Tr} A_S^2\le \|A_S\|\operatorname{Tr}A_S\). With \(S\) a uniform subset, independent of \(\pi\), we obtain
\begin{equation}
\begin{aligned}
Z_{2m}(q)
&\le Z_m(q)\,\mathbb E_{\pi\sim\nu_m}\prod_{\text{cycles }C_1\text{ of }\pi}
 \big(1+(q-1)2^{-|C_1|}\big)\\
&\le Z_m(q)\, Z_m(1+(q-1)/2).
\end{aligned}
\label{cas:e3}
\end{equation}
\end{proof}

\medskip\noindent\textbf{A slowly growing alphabet.}\enspace
We now derive the stated bound for $Z_b$ from the recursion.  This proof
retains the decay obtained by following the alphabet parameter until it
reaches its smallest value.
Let \(h_w(r)=2^{-w}\log Z_{2^w}(1+2^r)\) for integers \(r\ge 0\). For \(w\ge 1,r\ge 1\), \eqref{cas:e3} gives
\[
h_w(r)\le \tfrac12 h_{w-1}(r)+\tfrac12 h_{w-1}(r-1).
\]
At the boundary $r=0$, concavity gives
$Z_m(3/2)\le Z_m(2)^\alpha$, where $\alpha=\log_2(3/2)$.
Consequently
\[
 h_w(0)\le\theta h_{w-1}(0),\qquad
 \theta=(1+\alpha)/2<1,
 \qquad h_0(r)\le(r+1)\log2.
\]
To iterate these inequalities, start at $r$ and, while positive, stay
in place or decrement by one with equal probabilities. After reaching
zero, remain there and multiply by $\theta$ at each subsequent step.
The recurrence bounds $h_w(r)$ by $(r+1)\log2$ times the expected
product of these factors over $w$ steps.

Suppose $r\le w/8$. If $B$ counts successes in $\lfloor w/2\rfloor$
independent fair trials, the probability that the process has not
reached zero by that time is at most
\[
 \Pr(B<r)\le 2^r\E 2^{-B}
 =2^r(3/4)^{\lfloor w/2\rfloor}\le Ce^{-cw},
\]
since $2^{1/8}\sqrt{3/4}<1$. On the complementary event the product
of factors is at most $\theta^{w/2}$. Thus
\[
 h_w(r)\le(r+1)\log2\{Ce^{-cw}+\theta^{w/2}\}.
\]
Taking $r=\lfloor w/8\rfloor$ proves~\eqref{cas:e1}, with an absolute
$\gamma>0$ decreased as necessary to absorb the linear factor in $w$.

\medskip\noindent\textbf{Conditional cycle counts.}\enspace
The trace bound concerns an unconditioned network of one size.  To use it
on the many levels of a larger network, we leave a specified interval of
switch levels fresh and require a bound uniform in all other switches.
We next use \eqref{cas:e1} for switches sampled in the recursion defining \(\nu_n\). Nodes of size \(2^j\) will be called level \(j\) (\(1\le j\le d\)), with \(n/2^j\) nodes at that level indexed by fixed prefixes of length \(d-j\). Regard the input and output outer switches of every node as its switch variables at that level; the raw switch variables are independent. For a realization, draw at a node the edges taken by all its cards between its input pair indices and output pair indices (two parts); the graph is bipartite 2-regular, allowing double edges. Its two full matchings, together giving the edges with multiplicity, come from the two child bijections \(A,B\) on the \(m=2^{j-1}\) pair indices: the outer switches do not change these indices and there is one card per input pair going to each child. The number of cycle components, including cycles of two edges, is \(\kappa(A^{-1}B)\) (traverse using \(B\)-edges forward, \(A\)-edges backward). Let \(C_j\) be the sum of these numbers over level \(j\). It depends only on switches at levels strictly below \(j\), regardless of card identities coming into the nodes.

Fix \(1\le k\le n\), and put \(W_j=\min(k/2,C_j)\). With
\[
\delta=1/64
\]
we claim the absolute bound
\begin{equation}
\log\mathbb E\,2^{\delta\sum_{j=1}^d W_j}
 \le C k\log(e n/k).
\label{cas:e4}
\end{equation}
For \(j\ge 64\) let
\[
w=\lfloor j/2\rfloor,\qquad b=2^w,\qquad I_j=\{j-w,\ldots,j-1\}.
\]
Condition on all level switch variables outside \(I_j\). We will bound the exponential moment at the higher exponent \(\delta j\). Take \(q=1+2^{\lfloor w/8\rfloor}\), so \(q\ge 2^{\delta j}\) (\(\lfloor w/8\rfloor=\lfloor j/16\rfloor\ge j/64\) for \(j\ge64\)). Within each child of size \(m=2^{j-1}\) of a node at level \(j\), the interval \(I_j\) supplies \(w\) outer layers of input switches and the corresponding output layers. We can additionally condition on all the output switch variables there. The child input switches in \(I_j\) act as independent \(w\)-bit sweeps in \(m/b\) blocks (first \(w\) bits of the child position index used, its other bits fixed within a block); the blocks are aligned across the two children. Thus for a node, \(A^{-1}B\) now has the form
\[
U_A^{-1} G U_B
\]
where \(G\) is a fixed permutation of the \(m\) slots and \(U_A,U_B\) are independent permutations with law the product of \(\mu_b\) over those blocks. Indeed each child map first applies the indicated input switches, then its deeper fixed submaps and its conditioned output switches. In the \(q\)-color tensor representation on \(m\) slots, let \(T_{\rm bl}\) be the mean operator of the input sweeps in blocks, and \(U(G)\) the representation of \(G\). Conditionally,
\[
\mathbb E\,q^{\kappa(A^{-1}B)}
=\operatorname{Tr}\big(T_{\rm bl}^* U(G) T_{\rm bl}\big)
\le \operatorname{Tr}(T_{\rm bl}T_{\rm bl}^*)=Z_b(q)^{m/b}.
\]
Here we bounded the absolute trace against a positive semidefinite matrix using unitarity, and used the tensor product over blocks and cyclicity. Different level-\(j\) nodes use independent remaining randomness. The bound is uniform so we can drop the extra output conditioning. We get, with \(\mathcal F_{\overline I_j}\) denoting the switches outside \(I_j\),
\begin{equation}
\log\mathbb E\big[2^{\delta j W_j}\mid\mathcal F_{\overline I_j}\big]
\le \min\big( (\log 2)\delta j k/2,\ C n b^{-\gamma}\big).
\label{cas:e5}
\end{equation}
Indeed the second bound applies to \(C_j\) by multiplying over the \(n/2^j\) nodes and using \eqref{cas:e1} with \((n/2^j)(m/b)b^{1-\gamma}\le n b^{-\gamma}\).

\medskip\noindent\textbf{Combining overlapping level intervals.}\enspace
Estimate~\eqref{cas:e5} is conditional, so it does not require different
cycle counts to be independent.  We will use it first on collections of
disjoint intervals, and then average those collections by H\"older.
To combine intervals without paying a H\"older exponent of more than \(j\) on level \(j\), observe that for a fixed level \(l\),
\[
\sum_{j\ge 64:\ l\in I_j} \frac1j
\le \sum_{j=l+1}^{2l} \frac1j \le 1 .
\]
Using only \(64\le j\le d\), we can find weights \(\theta_{\mathcal S}\ge 0\), summing to 1, on subcollections \(\mathcal S\) of these indices with pairwise disjoint intervals, such that \(\sum_{\mathcal S\ni j}\theta_{\mathcal S}=1/j\) for each such \(j\). For clarity, assign to each interval, in order of increasing left endpoints (ties ordered arbitrarily), a measurable subset of \([0,1]\) of length \(1/j\), disjoint from the sets of earlier overlapping intervals. This is possible because those earlier overlapping intervals all contain the current left endpoint, so their sets' total length leaves enough room. Partition by membership in the assigned sets, taking lengths as weights. Weighted Hölder now gives
\[
\mathbb E\,2^{\delta\sum_{64\le j\le d} W_j}
\le \prod_{\mathcal S}\left(\mathbb E\,2^{\delta\sum_{j\in\mathcal S} j W_j}\right)^{\theta_{\mathcal S}}
\]
(positive weights suffice in the product, and the exponents inside the expectations on the right give the one on the left when averaged with those weights). Within a collection, use \eqref{cas:e5} on the largest \(j\) first and continue downwards. The terms from smaller \(j'\) in the collection depend only on levels \(\le j'-1\), hence outside \(I_j\) since \(I_{j'}\) ends at \(j'-1\) and the intervals are disjoint with \(j'-1<j-1\). Taking logs, the resulting bounds \eqref{cas:e5} occur with factors \(1/j\). Handling \(j<64\) deterministically, we obtain an upper bound for the left hand side of \eqref{cas:e4} by
\[
C k + \sum_{j\ge64} \min\big(Ck,\, Cn\,2^{-\gamma\lfloor j/2\rfloor}/j\big)
\le C' k\log(e n/k),
\]
where the sum bound can use all \(j\ge 64\): after \(j_0=\lceil (2/\gamma)\log_2(n/k)\rceil\), \(n2^{-\gamma\lfloor j/2\rfloor}\le 2^\gamma k\, 2^{-(\gamma/2)(j-j_0)}\). Thus the tail costs at most \(Ck\), and there are only \(O(\log(e n/k))\) terms before it. This proves \eqref{cas:e4}.

\medskip\noindent\textbf{From cycle weights to an injection density.}\enspace
The cycle moment is now available under the actual switch law.  The next
calculation relates it to the auxiliary uniform choices of legal child
colors.  Keeping these two probability laws separate is necessary when
we estimate a power of the density.
Denote by \([n]_{\ne}^k\) the ordered distinct \(k\)-tuples of positions, of cardinality \((n)_k=n(n-1)\cdots(n-k+1)\); positions can again be indexed by bit words, and \(\pi x\) denotes acting by the permutation \(\pi\) on every coordinate of \(x\). The following row bound is what we need:
\begin{equation}
\begin{gathered}
K_{n,k}(x,y)=\Pr_{\pi\sim\nu_n}(\pi x=y),\qquad
 R_x(y)=(n)_kK_{n,k}(x,y)\\
\Longrightarrow\quad
 \|R_x\|_{1+\delta}\le\exp(Ck\log(en/k)).
\end{gathered}
\label{cas:e6}
\end{equation}
for every such \(x\), with norm on uniform \(y\in[n]_{\ne}^k\).

Here is a partial route-counting identity. Fix the input and output tuples at a node. Draw the selected edges (the specified cards) between its input and output pair indices. Their components with edges are paths or cycles, say there are \(p_0\) paths and \(s\) cycles, with double edges forming a cycle as well. There are \(k+p_0\) vertices incident to edges (here \(k\) refers to the number of cards specified at this node). A proper edge coloring with two colors corresponds to assigning child halves, with two choices on each component by alternation. For each such coloring, the probability of the required outer switches, input and output together, is \(2^{-(k+p_0)}\): exactly one switch is prescribed per incident vertex, and the prescriptions are consistent since the tuples are distinct and the coloring is proper. In each child we now prescribe input and output tuples on the reduced indices, again distinct, for the tags of its color (keep their given order). With an empty tuple the kernel is interpreted as 1. Multiplying by \(n^k\) for a node of size \(n=2m\), the identity at that node is
\[
n^k K_{n,k}(x,y)
=2^s\,\mathbb E_{\rm color}
  \left[\prod_{a=0}^1 m^{k_a} K_{m,k_a}(x'_a,y'_a)\right]
\]
where the proper coloring is uniform and the child data (sizes \(k_a\)) are induced by it. At size 1 the normalized kernel is 1. Recursing defines a base coloring experiment for fixed \(x,y\), with independent uniform coloring choices at each level's nodes conditional on the already induced data there, inducing child data all the way down. Let \(S\) be the sum of selected-graph cycle counts \(s\) over its nodes. Thus \(n^k K_{n,k}=\mathbb E_{\rm base} 2^S\).

More precisely, for any outcome of this coloring tree, the probability under the actual full switches of taking \(x\) to \(y\) with those specified child assignments along the way is \(n^{-k}2^S\) times its probability in the base experiment. Indeed at each internal node \(v\), with local tuple size and component counts \(k_v,p_v,s_v\) (cards, paths with edges, cycles), the probability of the needed outer switches is \(2^{-(k_v+p_v)}\), versus \(2^{-(p_v+s_v)}\) for the designated coloring under the base experiment given its incoming data. For this fixed tree the switch requirements are prescriptions of raw coins, disjoint across the nodes, and together necessary and sufficient: tuples remain injective in each child (proper coloring) and the induced endpoints at size 1 require nothing more, with all unselected cards free of endpoint restrictions. The ratio thus multiplies to \(\prod_v 2^{s_v-k_v}=2^S n^{-k}\), since the \(k_v\)'s sum to the original \(k\) on each level. Conversely a full switch realization from \(x\) determines a unique selected endpoint and coloring tree for it (actual child routes).

On a realization of the full switches starting with tuple \(x\), call the corresponding total cycle count \(S_x\), using its actual endpoint and child routes. Any cycle in a selected partial graph is a full cycle component of that node's graph from the realization, since its vertices already have degree 2 on the cycle. Also each uses at least two selected cards, with disjoint cards across cycles of a level. It follows that \(S_x\le\sum_j W_j\) (now with the original full tuple size \(k\)) on every realization. By convexity and the refined identity, with \(u=(n)_k/n^k\le 1\),
\[
\mathbb E_{y\ {\rm uniform}} R_x(y)^{1+\delta}
\le \frac{u^{1+\delta}}{(n)_k}\sum_y \mathbb E_{\rm base(x,y)} 2^{(1+\delta)S}
= u^\delta \,\mathbb E_{\rm actual} 2^{\delta S_x}.
\]
Now \eqref{cas:e4} proves \eqref{cas:e6}.

\subsection{A contraction measured by the transposition Casimir}

The density estimate is complete.  We now prove the operator contraction
on the same scale $h_n(k)$.  The induction splits the positions into two
halves, so its main task is to compare the parent Casimir scalar with the
two child scalars seen by a vector fixed by the outer switches.

We use standard representation theory of the symmetric group over \(\mathbb C\): irreducibles are indexed by partitions \(\lambda\) of \(n\), of dimensions \(f^\lambda\) (numbers of standard Young tableaux). We use Young's branching rule for restriction down the symmetric groups (multiplicity given by paths of successively removing single boxes while keeping diagrams of partitions), and the standard transposition-sum content formula (by the Young--Jucys--Murphy content theorem~\cite[Eq.~(2.1), Proposition~5.3 and Theorem~5.8]{VershikOkounkov2005}), that \(\sum_{i<j}(i\,j)\) acts in \(\lambda\) by the sum of column-minus-row indices of its boxes. Indices of rows and columns here start at 1. The permutation module on ordered distinct \(k\)-tuples is the coset module with stabilizer \(S_{n-k}\), so multiplicity of \(\lambda\) is the dimension of its stabilizer-fixed space by Frobenius reciprocity. In particular for \(k=n-\lambda_1\), \(\lambda\) occurs in this module and (if \(k\ge1\)) not in the corresponding \(k-1\) tuple module. Its multiplicity at this \(k\) is \(f^{\bar\lambda}\), where \(\bar\lambda\) consists of all rows below the first: paths in the branching rule to the trivial shape \((n-k)\) must remove just that tail. Also
\begin{equation}
f^\lambda\le \binom{n}{k} f^{\bar\lambda}
\label{cas:e7}
\end{equation}
by choosing the entries of a standard tableau below the top row.

Let $U$ denote the unitary action of $S_n$ on the representation currently
under consideration.  Write \(D_{ij}=I-U((i\,j))\) on a unitary representation (indices \(i,j\) are positions). They are positive semidefinite. Put \(\mathcal C=\sum_{i<j}D_{ij}\), which in \(\lambda\) acts by
\begin{equation}
c=c_\lambda=\tfrac12(n^2-\textstyle\sum_i\lambda_i^2)+\sum_i(i-1)\lambda_i,\qquad
nk/2\le c\le nk,\qquad k=n-\lambda_1.
\label{cas:e8}
\end{equation}
For the lower bound use \(\sum\lambda_i^2\le n\lambda_1\). For the upper, the sum of squares below row 1 is at least \(k\), and the sum of row-minus-one indices over tail boxes is at most \(k(k+1)/2\) (ordering tail boxes by row, the \(a\)-th has row index at most \(a+1\)). Define
\[
\Phi_n(c)=\frac{c}{n}\big(2+\log(n^2/c)\big),\qquad \Phi_n(0)=0.
\]
We prove there are absolute \(B>1,\eta>0\) such that on every nontrivial irreducible \(\lambda\) for \(n=2^d\),
\begin{equation}
\|K_n\|_{\lambda}\le B^{-1}\exp(-\eta\Phi_n(c_\lambda)).
\label{cas:e9}
\end{equation}
The norm notation means the operator norm in the indicated irreducible. In particular the bound implies \(\|K_n\|_\lambda\le\exp(-a k\log(e n/k))\) there for an absolute \(a>0\): \(k/2\le c/n\le k\) gives \(\Phi_n(c)\ge(k/2)(2+\log(n/k))\).

We will do an induction for \eqref{cas:e9} after some random-split estimates. In outline, \(K_n\) is the child product-law operator compressed by \(\Pi_1\), so we study Rayleigh values in unit vectors invariant under the bit-1 switch group. Under restriction to the two halves' product subgroup, the induction hypothesis will contribute two child \(\Phi\)-values (evaluated at the partial Casimir scalars on the irreducibles of the factors), in place of the \(\Phi_n(c)\) in \eqref{cas:e9}. The invariant vector lets us analyze the distribution of the partial Casimir values (spectral weights in the vector) by randomly flipping which member of each pair lies in each half. We control the sum of these two scalars relative to \(c/2\), and their difference, for the \(\Phi\)-comparison; we also use the high probability of two nontrivial child factors in the spectral weighting for large \(k\) so that the factor \(B^{-1}\) in \eqref{cas:e9} gives slack.

\subsubsection{Two estimates for a random split}

Partition the positions into the \(m=n/2\) switch pairs of bit 1, using subscripts \(p0,p1\) in a pair \(p\) for its bit-1 values 0 and 1, respectively. Draw independent fair signs \(s_p\). Write \(\xi_{p0}=s_p,\ \xi_{p1}=-s_p\), and let \(L,R\) be the positions with \(\xi=+1,-1\) respectively. The two partial Casimirs \(\mathcal C_L,\mathcal C_R\) are the sums of \(D_{ij}\) for pairs of positions both in the indicated side. They commute for a fixed split. Use the operators
\[
\Delta=\mathcal C-2(\mathcal C_L+\mathcal C_R),\qquad
\mathcal R=\mathcal C_L-\mathcal C_R.
\]
On \(\lambda\), both \(\Delta\) and \(\mathcal R\) are between \(-cI\) and \(cI\) by positivity of the terms of \(\mathcal C\). The sum defect $\Delta$ measures how much Casimir mass is lost in the
split, while $\mathcal R$ measures the difference between the two sides.
They require different estimates: an exponential moment of $\Delta$
itself and an exponential moment of $\mathcal R^2$.

\begin{lemma}[Exponential moments for a transversal split]\label{cas:split-moments}
There are absolute constants $t_0,t_1>0$ and $C_0<\infty$ such that,
for every dyadic $n\ge2$ and every nontrivial $\lambda\vdash n$, with
$k=n-\lambda_1$ and $c=c_\lambda$,
\begin{equation}
\left\|\mathbb E_s e^{t\Delta/n}\right\|_\lambda\le C_0
\quad\text{at }t=t_0\log(e n/k),\qquad
\left\|\mathbb E_s e^{t_1\mathcal R^2/(nc)}\right\|_\lambda\le C_0,
\label{cas:e10}
\end{equation}
The expectations are over the independent fair signs defining the
transversal split, and both norms are operator norms on $V_\lambda$.
\end{lemma}
\begin{proof}
\medskip\noindent\textbf{The imbalance moment.}\enspace
We start with the second estimate. Write \(F_i=\sum_{j\ne i} D_{ij}\). Then \(\mathcal R=\tfrac12\sum_i\xi_i F_i=\sum_p s_p W_p\) with \(W_p=(F_{p0}-F_{p1})/2\). On \(\lambda\),
\[
K_0=\max_p\|W_p\|\le n,\qquad M^2=\left\|\sum_p W_p^2\right\|\le 2nc,
\]
using \(0\le F_i\le 2(n-1)I\), \(F_i^2\le2(n-1)F_i\), \(\sum_i F_i=2cI\), and \(\sum_p W_p^2\le \tfrac12\sum_i F_i^2\) (by \((F_{p0}+F_{p1})^2\ge0\) for each pair). For such self-adjoint matrices the following moment bound suffices, for integers \(l\ge1\):
\begin{equation}
\left\|\mathbb E_s(\textstyle\sum s_p W_p)^{2l}\right\|
\le \big(C(M\sqrt l+K_0 l)\big)^{2l}.
\label{cas:e11}
\end{equation}
Here are details without a dimension factor. Expand the noncommuting product. In the expectation of the scalar signs, for sign indices \(p_1,\ldots,p_{2l}\) we have
\[
\mathbb E_s\prod_{u=1}^{2l} s_{p_u}
=\sum_{\mathcal B}\prod_{B_0\in\mathcal B}
   \left(\mathbf1_{\{p_u\text{ all equal for }u\in B_0\}}\,\chi_{|B_0|}\right),
\qquad \chi_b=\left.\frac{d^b}{dz^b}\log\cosh z\right|_{z=0}
\]
over set partitions \(\mathcal B\) of the \(2l\) places. Indeed the log of the generating expectation for \(\exp(\sum_u z_u s_{p_u})\) is \(\sum_p\log\cosh(\sum_{u:\,p_u=p} z_u)\) as power series near 0. Exponentiating and taking the multilinear coefficient gives the expansion. Singletons have zero coefficient, and \(|\chi_b|\le b! C^b\) by analyticity near 0.

In the term for a fixed partition with \(h\) blocks of sizes at least 2, there are within-block index equalities in the sum of operator words \(W_{p_1}\cdots W_{p_{2l}}\), but indices across blocks can be summed without distinctness restrictions. The norm of this contracted sum (before the factors \(\chi_b\)) is bounded by \(M^{2h}K_0^{2l-2h}\). To see this, apply from right to left, keeping an extra index register \(\mathbb C^m\) (with orthonormal basis indexed by pairs \(p\)) per open block in tensor product with the representation. On first encountering a block, use the column map \(x\mapsto\sum_p W_p x\otimes e_p\) of norm \(M\), inserting its register; on final occurrence use its adjoint (\(W_p=W_p^*\)); on intermediate occurrences use \(\sum_p W_p\otimes |e_p\rangle\langle e_p|\) with that register, of norm at most \(K_0\). At each step tensor with identities on other open registers and permute registers as needed. This composition realizes precisely the sum over assignments of indices to blocks, as the first encounter branches over the new index, intermediate encounters use that same index, and the final one sums it out.

Summing absolute bounds using ordered compositions \(b_1+\cdots+b_h=2l\) with \(b_i\ge2\), partitions are accounted for with factor \((2l)!/(h!\prod_i b_i!)\). The \(b_i!\) cancel in the bound from \(\prod_i(b_i! C^{b_i})\), there are at most \(2^{2l}\) compositions in total, and \((2l)!/h!\le (2l)^{2l-h}\). Now \(l^{2l-h}M^{2h}K_0^{2l-2h}=(M\sqrt l)^{2h}(K_0 l)^{2l-2h}\) with \(h\le l\), giving \eqref{cas:e11}.

For \(1\le l\le k\), \eqref{cas:e11} and \eqref{cas:e8} imply \(\|\mathbb E_s(\mathcal R^2/(nc))^l\|_\lambda\le (C l)^l\) (\(M/\sqrt{nc}\le\sqrt2,\ K_0/\sqrt{nc}\le\sqrt{2/k}\)). For \(l>k\), use \(\|\mathcal R^2/(nc)\|\le c/n\le k\) on \(\lambda\) instead, so the same bound holds. The exponential series, using \(l!\ge(l/e)^l\), now proves the \(\mathcal R\) bound in \eqref{cas:e10} at a sufficiently small \(t_1>0\).

\medskip\noindent\textbf{The sum defect on injections.}\enspace
The imbalance bound is now proved.  For the sum defect we will first
replace signs by an entrywise majorant, then bound that majorant through
a Gaussian tensor calculation.  This larger carrier avoids a factor
$D_\lambda$ in the estimate.
We prove the first estimate in~\eqref{cas:e10} on the whole ordered
distinct $k$-tuple module $\mathcal V=\ell^2(X_{n,k})$, which contains
$V_\lambda$. Embed $\mathcal V$ in $(\mathbb C^n)^{\otimes k}$ with
the standard ordered tuple basis, so $U(g)|x\rangle=|gx\rangle$, and
write $D_{\ne}$ for its orthogonal projection. Let
$\mathcal V_{\rm sym}\subseteq\mathcal V$ consist of vectors invariant
under permutations of the tensor slots. For a real vector $u$ on
positions put $p_u=|u\rangle\langle u|/n$ and
$\mathcal L^u=\sum_{a=1}^k p_u^{(a)}$. On $\mathcal V$, the identity
$\sum_i\xi_i=0$ gives
\begin{equation}
\Delta=-\sum_{i<j}\xi_i\xi_j D_{ij}
=n D_{\ne}\mathcal L^\xi D_{\ne}+ B^\xi,\qquad
(B^\xi f)(x)=\sum_{a<b}\xi_{x_a}\xi_{x_b}\big(f(x)+f(x^{a\leftrightarrow b})\big),
\label{cas:e12}
\end{equation}
where the tuple on the right interchanges two slots. In fact in the transposition sum one separates replacements of a single occupied position by an unoccupied one from interchanges of two occupied positions. Both have off-diagonal coefficient \(\xi_i\xi_j\) for positions \(i,j\) involved. The diagonal at \(x\), for its occupied set \(X\), equals \(-\sum_{i\in X,\,j\notin X}\xi_i\xi_j-\sum_{i<j:\,i,j\in X}\xi_i\xi_j=k+\sum_{i<j:\,i,j\in X}\xi_i\xi_j\).

\medskip\noindent\textbf{Removing the opposite signs inside a switch pair.}\enspace
All entries of powers of the last expression in \eqref{cas:e12} on distinct tuples are given by polynomials in the position variables with nonnegative coefficients before substituting the signs. Replace \(\xi\) by \(v\) with \(v_{p0}=v_{p1}=s_p\) in this expression (the resulting operator \(\widetilde\Delta=nD_{\ne}\mathcal L^vD_{\ne}+B^v\) need not equal the transposition formula with that replacement). For any such monomial with the original signs, the absolute value of its average is bounded by its average with the replacement: both averages vanish on odd pair-index parity at any \(p\), and the replacement gives \(+1\) on all even parities. For \(t\ge0\), it follows by the series and triangle inequality that entrywise
\[
\left|(\mathbb E_s e^{t\Delta/n})_{xy}\right|
\le (\mathbb E_s e^{t\widetilde\Delta/n})_{xy}.
\]
The latter matrix is real symmetric entrywise nonnegative and commutes with permutations of the tensor slots (in fact \(\widetilde\Delta\) commutes with them for every set of signs). Its operator norm, which bounds the former norm, is attained on slot-symmetric distinct tensors. Indeed a symmetric entrywise nonnegative matrix has a nonnegative eigenvector for its largest eigenvalue realizing the spectral radius (or maximize using entrywise absolute values in the Rayleigh bound), and here it can be averaged over slot permutations without vanishing. Entrywise domination in absolute value by nonnegative matrices bounds operator norm by applying them to entrywise absolute vectors.

On the slot-symmetric distinct subspace use the orthonormal basis indexed by \(k\)-sets \(X\), the normalized sums of orderings. In this basis \(\widetilde\Delta\) has diagonal \(V_s(X)=(\sum_{i\in X} v_i)^2\), coming from \(k\) on the diagonal of \(nD_{\ne}\mathcal L^vD_{\ne}\) and \(2\sum_{i<j:\,i,j\in X}v_i v_j\) from \(B^v\) on symmetric functions. Its off-diagonal is \(v_i v_j\) between \(X\) and \((X\setminus\{i\})\cup\{j\}\) (\(i\in X,j\notin X\)), zero otherwise. Let \(J\) be the unweighted adjacency matrix of this single-exchange graph. Conjugating \(\widetilde\Delta\) by the diagonal signs \(v_X=\prod_{i\in X} v_i\) gives \(J+V_s\) with \(V_s\) diagonal. Thus in computing \((e^{t\widetilde\Delta/n})_{XY}\) by this conjugation, the sign factor is \(v_Xv_Y=\prod_{p\in O}s_p\), where \(O=O(X,Y)\) consists of pair indices of odd combined occupancy in \(X,Y\) (counting multiplicity). Write \(h=|O|\), always even.

\medskip\noindent\textbf{A Gaussian bound along each path.}\enspace
We need to average the endpoint sign $v_Xv_Y$ together with the
nonnegative potential $V_s$.  An expansion over paths keeps this sign
visible and turns the potential into a Gaussian covariance.
In the exponential kernel of \((J+V_s)/n\) over duration \(t\), expand in graph paths using \(J/n\), with a multiplier on each timed path \(X(u)\) given by
\[
\exp\left(\frac1n\int_0^t V_s(X(u))\,du\right).
\]
Specifically, using symmetry of the kernel, its \(X,Y\) entry is
\[
\sum_{r=0}^\infty n^{-r}
\sum_{X=X_0\sim X_1\sim\cdots\sim X_r=Y}
\int_{0<u_1<\cdots<u_r<t}
\exp\left(n^{-1}\sum_{i=0}^r(u_{i+1}-u_i)V_s(X_i)\right)\,du_1\cdots du_r ,
\]
where \(\sim\) denotes adjacency in \(J\), \(u_0=0,\ u_{r+1}=t\), \(X(u)=X_i\) between successive times \(u_i,u_{i+1}\), and \(r=0\) means no jump (\(X=Y\)). This is the ordinary iterated exponential expansion (Duhamel iteration with the exponential of the diagonal term retained between jumps); it converges absolutely by boundedness in finite dimension and the simplex volumes \(t^r/r!\). The nonnegative baseline weights without the multipliers sum to \((e^{tJ/n})_{XY}\). On a fixed timed path let \(b_p(u)\in\{0,1,2\}\) denote occupancy at pair \(p\), so \(\sum_p b_p=k\). Define a centered real Gaussian \(z\) on the \(m\) pair indices with covariance
\[
H=\frac2n\int_0^t b(u)b(u)^\top\,du .
\]
The path multiplier is \(\mathbb E_z\exp(\sum_p z_p s_p)\). Thus after including \(v_Xv_Y\), averaging the signs on the path gives \(\mathbb E_z\prod_{p\in O}\sinh z_p\prod_{p\notin O}\cosh z_p\), bounded in absolute value by \(\mathbb E_z e^{\|z\|_2^2/2}\prod_{p\in O}|z_p|\), using \(\cosh y\le e^{y^2/2}\), \(|\sinh y|\le |y| e^{y^2/2}\) (by the series). We have \(\operatorname{Tr}H\le 4kt/n\), \(H_{pp}\le8t/n\). From now on take \(t=t_0\log(e n/k)\) for an absolute \(t_0>0\) sufficiently small, starting with \(t_0\le1/8\); then \(\operatorname{Tr}H\le1/2\) since \((k/n)\log(e n/k)\le1\). Tilting the Gaussian by \(e^{\|z\|_2^2/2}\) gives mass \(\det(I-H)^{-1/2}\le e^{\operatorname{Tr}H}\) and covariance \(H(I-H)^{-1}\le2H\) after normalization, as seen in an eigenbasis (allowing zero eigenvalues, and using \(-\log(1-y)\le2y\) for \(0\le y\le1/2\)). By scalar Gaussian absolute moments and Hölder on the \(h\) factors, the sign-averaged path multiplier with the endpoint sign included has absolute value at most
\[
C\big(C_1\sqrt{t h/n}\big)^h
\]
with the power taken as 1 if \(h=0\). The constants \(C,C_1\) here can be taken absolute for all \(0<t_0\le1/8\), using e.g. \(\mathbb E|Z|^h=(h-1)!!\le h^{h/2}\) for a standard normal \(Z\) and positive even \(h\). Uniformly summing the path weights, we have on this symmetric subspace the entrywise nonnegative averaged matrix \(\mathbb E_s e^{t\widetilde\Delta/n}\) bounded entrywise above by
\begin{equation}
C(e^{tJ/n})_{XY} \big(C_1\sqrt{t h/n}\big)^h.
\label{cas:e13}
\end{equation}

\medskip\noindent\textbf{Absorbing the endpoint factor into a tensor operator.}\enspace
The path estimate~\eqref{cas:e13} still depends on $h$, the number of pair
indices with odd endpoint occupancy.  A common Gaussian coordinate will
supply exactly the moments needed to absorb this factor.  After that
comparison, only the norm of one explicit positive tensor average remains.
We absorb this extra weight into a tensor operator. With \(\mathbf 1\) the all-ones vector on positions, on the same symmetric distinct basis \((J+kI)/n\) is the restriction of \(D_{\ne}\mathcal L^{\mathbf 1} D_{\ne}\). All slot sums here commute with slot permutations. We can drop a scalar factor \(e^{-tk/n}\), and in the remaining exponential power series use entrywise nonnegativity in the ordered basis to drop the distinctness requirement on intermediates (then compare on normalized sums of orderings). Thus \(e^{tJ/n}\) is bounded entrywise by the symmetric-distinct compression of
\[
e^{t\mathcal L^{\mathbf 1}}=\bigotimes_{a=1}^k (I+\beta p_{\mathbf 1}),\qquad \beta=e^t-1.
\]
To dominate the weight depending on \(O(X,Y)\) as well, take a duplicated Gaussian vector \(g\) on positions with
\[
g_{p0}=g_{p1}=g'_p+y_*, 
\]
where \(g'_p\) are independent standard real Gaussians and \(y_*\) is an independent centered real Gaussian of variance \(C_* t/n\). Then \eqref{cas:e13} is bounded entrywise by \(C\) times the symmetric distinct compression of
\begin{equation}
G_t:=\mathbb E_g \bigotimes_{a=1}^k (I+\beta p_g).
\label{cas:e14}
\end{equation}
Indeed expand each tensor product, checking on ordered endpoints \(x,y\) with sets \(X,Y\). Identity slots (inactive in a contributing term) require the same position at both ends. Thus in the active slots \(E\subseteq [k]\), the occurrences of pair indices among \(x_a,y_a\) for \(a\in E\) have the same parities (odd-index set \(O(X,Y)\)). In \eqref{cas:e14} we get the additional factor \(\mathbb E_g\prod_{a\in E}g_{x_a}g_{y_a}\) compared to the tensor expansion with \(p_{\mathbf 1}\). The Gaussian monomial expectations are nonnegative; expanding \(g\), assign for each odd pair-index one occurrence to \(y_*\) and the remaining occurrences everywhere to the independent \(g'_p\), all the latter degrees then even. Those even moments of standard normals are at least 1, and for positive even \(h\),
\(\mathbb E y_*^h=(C_*t/n)^{h/2}(h-1)!!\ge (C_1\sqrt{th/n})^h\)
by taking \(C_*\) absolutely large enough; for instance \((h-1)!!=h!/(2^{h/2}(h/2)!)\ge (h/e^2)^{h/2}\). This also works with moment 1 at \(h=0\). Other assignments contribute nonnegative expectations. Summing over orderings with the basis normalizations for the symmetric-distinct compression preserves the entrywise bound.

Let $Q:\mathcal V_{\rm sym}\longrightarrow(\mathbb C^n)^{\otimes k}$
be the isometric inclusion, whose $k$-set basis vectors are normalized
sums of their orderings. The comparisons just proved give
\begin{equation}\label{cas:gaussian-norm-chain}
\begin{aligned}
 \|\E_s e^{t\Delta/n}\|_{V_\lambda}
 &\le\|\E_s e^{t\Delta/n}\|_{\mathcal V}\\
 &\le\|\E_s e^{t\widetilde\Delta/n}|_{\mathcal V_{\rm sym}}\|_{\op}\\
 &\le C\|Q^*G_tQ\|_{\op}\le C\|G_t\|_{\op}.
\end{aligned}
\end{equation}
The middle inequalities use entrywise domination in the indicated
orthonormal bases and the nonnegative eigenvector argument above.
The last inequality is the norm bound for an orthogonal compression.
It remains to prove $\|G_t\|_{\op}\le2$ on the full tensor product.

\medskip\noindent\textbf{Bounding the Gaussian tensor average.}\enspace
On the one-slot space, $p_g$ is supported on the pair-sum subspace
with orthonormal basis $(e_{p0}+e_{p1})/\sqrt2$. There it becomes
$|\widetilde g\rangle\langle\widetilde g|/m$, where
$\widetilde g_p=g'_p+y_*$ has covariance
$I+(C_*t/n)\mathbf1\mathbf1^\top$ in $m$ dimensions. Splitting each
slot into this subspace and its orthogonal complement makes $G_t$
block diagonal. Complementary slots contribute identity, so a block
with $r\le k$ pair-sum slots requires a bound only on
$(\mathbb C^m)^{\otimes r}$.

Diagonalize the covariance by a real orthogonal change of basis.
The new coordinates are independent centered Gaussians with variances
at most $\Gamma=1+C_*t/2$. In this basis, every averaged monomial
in the tensor expansion is zero or a nonnegative product of even
moments. Increasing all these variances to $\Gamma$ therefore gives
an entrywise upper bound, and hence an operator-norm upper bound.

The resulting comparison matrix for independent coordinates of variance
$\Gamma$ is entrywise nonnegative and commutes with slot permutations.
A nonnegative norm-maximizing eigenvector can again be averaged over
these permutations without vanishing. It therefore suffices to bound
the matrix on $\operatorname{Sym}^r(\mathbb C^m)$.

We have reduced the problem to symmetric tensors in $m$ independent
Gaussian coordinates.  The next calculation bounds each coefficient of
the tensor product.  Its dependence on the degree is what permits a
geometric sum even when $t$ grows like $\log(en/k)$.
Here are details of the Gaussian norm estimate there, now using \(\widetilde g\) as the vector for the i.i.d. coordinate comparison. Use creation and annihilation operators \(a_i^\dagger,a_i\) on the algebraic direct sum of symmetric tensor powers (only acting between finite degrees here), where on normalized occupation vectors with \(n_i\) copies of basis vector \(i\) they raise or lower \(n_i\) with coefficients \(\sqrt{n_i+1}\) or \(\sqrt{n_i}\), respectively (zero for lowering at 0). For the real vector \(\widetilde g\) write \(a(\widetilde g)=\sum_i \widetilde g_i a_i\). Restricted to symmetric degree \(r\), for \(0\le j\le r\) the sum over all size-\(j\) subsets of slots of the product of \(|\widetilde g\rangle\langle\widetilde g|/m\) in active slots is
\[
\frac{(a(\widetilde g)^\dagger)^j a(\widetilde g)^j}{j!\,m^j}.
\]
One can see this in an orthonormal basis with first vector along \(\widetilde g\), counting the choices of \(j\) copies among those in its first mode: on an occupation vector for this basis with first-mode occupancy \(n'_1\), both sides multiply by \(\binom{n'_1}{j}(\|\widetilde g\|_2^2/m)^j\). Here creation along a vector, from degree \(u\), equivalently inserts that vector then projects to symmetric tensors with the factor \(\sqrt{u+1}\), with lowering the adjoint, so creation/annihilation along a unit direction obey the same formulas in a basis containing that unit vector. (At the zero vector the identity for the subset sum is immediate.) Take expectation with the i.i.d. Gaussian coordinates. The scalar Gaussian pairing identity~\cite{Isserlis1918} gives an order-\(2j\) scalar moment as a sum over all pairings of the \(2j\) places, with factor \(\Gamma\mathbf1_{\{i=i'\}}\) per pair of coordinate indices \(i,i'\) (also seen by differentiating the Gaussian moment generating function). Applied to the numerator in the operator formula, it sums over pairings of the scalar coefficients of creators and annihilators, leaving normal order of the operators. Put \(A=\sum_i a_i^2\) and let \(\mathcal N=\sum_i a_i^\dagger a_i\) be tensor degree. For \(\ell\) creator-creator pairs there are also \(\ell\) annihilator-annihilator pairs and a bijective pairing across the remaining \(j-2\ell\) of each. Each fixed pairing in the numerator gives factor \(\Gamma^j\) times the contracted operator
\[
(A^\dagger)^\ell\,\mathcal N^{\underline{j-2\ell}}\, A^\ell,
\]
using falling powers \(x^{\underline{q'}}=x(x-1)\cdots(x-q'+1)\), including \(x^{\underline{0}}=1\). Indeed creators commute among themselves, as do annihilators among themselves. The cross terms in the middle sum a normally ordered removal and replacement of an ordered selection of \(q'=j-2\ell\) particles: on occupations \((n_i)\) this counts \((\sum_i n_i)^{\underline{q'}}\), since a given ordered list of modes with multiplicities \(d_i\) contributes \(\prod_i n_i^{\underline{d_i}}\). There are \(j!/((j-2\ell)!\,2^\ell\ell!)\) choices of the internal pairs on each side and \((j-2\ell)!\) bijections of the remaining places. Thus after division by \(j!\) the count coefficient is \(j!/((j-2\ell)!\,2^{2\ell}(\ell!)^2)\le \binom{j}{2\ell}\).

On degree \(u\), \(\|A\|^2\le u(u+m)\), writing the norm for the action out of that degree. This follows using \([a_i,a_j^\dagger]=\mathbf 1_{\{i=j\}}I\), hence \([A,A^\dagger]=4\mathcal N+2mI\). For \(u\ge2\) the squared norm out of \(u\) is the norm of \(A A^\dagger\) on \(u-2\), which equals the squared norm out of \(u-2\) plus \(4(u-2)+2m\); the sequence starts with 0 on degrees 0,1. This gives the bound by induction. For \(1\le j\le r\), out of degree \(r\) the squared norm of \(A^\ell\) is thus at most \((r(r+m))^\ell\) for \(2\ell\le j\), and the middle falling power of \(\mathcal N\) in the contracted operator is evaluated on degree \(r-2\ell\), bounded there by \(r^{j-2\ell}\). Consequently the norm of the averaged size-\(j\) sum on the symmetric subspace in this range is at most
\[
(\Gamma/m)^j\sum_{0\le \ell\le j/2}\binom{j}{2\ell} r^{j-2\ell}(r(r+m))^\ell
\le \left(\frac{\Gamma}{m}(r+\sqrt{r(r+m)})\right)^j
\]
(and use the bound 1 at \(j=0\) for the identity).

For sufficiently small absolute \(t_0\) this sums in \eqref{cas:e14}, since
\[
(e^t-1)(1+C_*t/2)\,\frac{k+\sqrt{k(k+m)}}{m}\le\frac12 .
\]
In fact the last ratio is at most $C\sqrt{x}$ for $x=k/n\in(0,1]$.
Since $e^t-1\le te^t$ and $t=t_0\log(e/x)$, the expression is bounded
by an absolute multiple of
\[
 t_0\log(e/x)x^{1/2-t_0}
 +C_*t_0^2\log(e/x)^2x^{1/2-t_0}.
\]
For $0<t_0\le1/8$, both powers of the logarithm times
$x^{1/2-t_0}\le x^{3/8}$ are uniformly bounded. With $C_*$ already
fixed by the endpoint comparison, we may therefore decrease $t_0$
until the displayed geometric ratio is at most $1/2$.

The sum with coefficients $\beta^j$, $0\le j\le r$, then has norm
at most two in every sector. Hence $\|G_t\|_{\op}\le2$.
Equation~\eqref{cas:gaussian-norm-chain} now bounds
$\|\E_s e^{t\Delta/n}\|_{V_\lambda}$ by an absolute constant,
completing~\eqref{cas:e10}.
\end{proof}

\subsubsection{Induction of the operator bound}

The two split moments now control the change in the potential
$\Phi_n$.  We first derive that scalar comparison, then use the extra
factor $B^{-1}$ from two nontrivial children to close the induction.
For the split with fixed halves (the children in the recursion), denote by \(H_1\) the outer switch group. In \(\lambda\ne(n)\), to bound the norm of \(K_n\) it suffices to bound its Rayleigh values on unit vectors \(v\) invariant under \(H_1\), since \(K_n\) is positive semidefinite and compressed by the projection \(\Pi_1\) onto these invariants (if there is no such vector the norm is zero). For such a vector, the Rayleigh value is that for the operator of the middle product law from the two independent children, since \(\Pi_1 v=v\). Under restriction to \(S_m\times S_m\), decompose into an orthogonal sum of product irreducibles (using the usual tensor-product form of irreducibles of a direct product of finite groups), with the joint partial Casimir values \(c_0,c_1\) from the two factors. We use \(c_0,c_1\) as classical random variables with the joint spectral probabilities in \(v\), i.e. squared norms of the joint spectral projections applied to \(v\). Their law for this fixed split is also their law for each split in the sign randomization used in \eqref{cas:e10}: any such transversal split is the image of the fixed one by an element of \(H_1\), which fixes \(v\), and the operators for the splits correspond by conjugation. Write
\[
D_0=c-2(c_0+c_1),\qquad D_1=c_0-c_1,\qquad
\mathcal E=\Phi_n(c)-\Phi_m(c_0)-\Phi_m(c_1).
\]
In this scalar law we have \(\mathbb E e^{tD_0/n}\le C_0\) and \(\mathbb E e^{t_1 D_1^2/(nc)}\le C_0\) by \eqref{cas:e10} evaluated against \(v\). Also \(|D_0|,|D_1|\le c\). Put \(u=c/n\) and \(x_i=4 c_i/c\) for \(i=0,1\), with mean \(\bar x\) of the two \(x_i\). Then
\[
\mathcal E=u\big[(1-\bar x)(2+\log(n/u))+\tfrac12\textstyle\sum_{i=0}^1 x_i\log x_i\big]
\le \frac{D_0}{n}(1+\log(n/u))+\frac{D_0^2+4D_1^2}{n c},
\]
where \(0\log0=0\), \(x\log x\le (x-1)+(x-1)^2\) for \(x\ge0\), and \(x_0-1=-D_0/c+2D_1/c,\ x_1-1=-D_0/c-2D_1/c\). Using \(|D_0|\le c\) and \(\log(n/u)\ge0\), the \(D_0\) terms on the right are at most \((D_0)_+(2+\log(n/u))/n\), writing \(z_+=\max(z,0)\). Here \(k/2\le u\le k\), so \(2+\log(n/u)\le C\log(e n/k)\). For an absolute \(a_0>0\) sufficiently small we have \(2a_0(2+\log(n/u))\le t,\ 8a_0\le t_1\). Cauchy--Schwarz under the scalar law therefore gives
\[
\mathbb E e^{a_0\mathcal E_+}
\le (\mathbb E e^{t(D_0)_+/n})^{1/2}
     (\mathbb E e^{t_1 D_1^2/(nc)})^{1/2}\le C
\]
using \eqref{cas:e10} as above and \(e^{t(D_0)_+/n}\le1+e^{tD_0/n}\). In particular
\begin{equation}
\mathbb E e^{\eta\mathcal E_+}\le 1+C'\eta\qquad (0\le\eta\le a_0)
\label{cas:e15}
\end{equation}
by decreasing the exponent via Hölder (\(C^{\eta/a_0}\le1+(C-1)\eta/a_0\) for \(C\ge1\)). Also
\begin{equation}
\Pr(c_0=0\text{ or }c_1=0)\le
\Pr(D_0\ge c/2)+\Pr(|D_1|\ge c/4)
\le C \exp(-c' k)
\label{cas:e16}
\end{equation}
for an absolute \(c'>0\). Indeed if one \(c_i=0\), either their sum is at most \(c/4\) or \(|c_0-c_1|\ge c/4\). The two probability terms are bounded by \(C_0 e^{-t c/(2n)}\) and \(C_0 e^{-t_1 c/(16n)}\) using \eqref{cas:e10}, with \eqref{cas:e8}. Take an absolute integer \(K\) large enough that for \(k\ge K\) the probability of a zero \(c_i\) in \eqref{cas:e16} is at most \(1/4\).

Assuming \eqref{cas:e9} on the nontrivial child irreducibles, the Rayleigh value above is bounded by
\[
\mathbb E\, B^{-z}\exp\big[-\eta(\Phi_m(c_0)+\Phi_m(c_1))\big],
\quad z=\mathbf 1_{\{c_0>0\}}+\mathbf 1_{\{c_1>0\}},
\]
by taking the operator norm product on each product summand (there the middle operator is the tensor product of the two child \(K_m\)'s in the respective irreducibles; trivial factors have norm 1 and \(c_i=0\)). For \(k\ge K\) and \(1<B\le2\),
\[
\mathbb E B^{1-z}\le 1-\tfrac18(B-1)
\]
by \eqref{cas:e16}, since \(z=2\) has probability at least \(3/4\) with deficit \(1-1/B\), and otherwise the excess over 1 is at most \(B-1\). Using \eqref{cas:e15} and \(B^{1-z}\le2\) gives
\[
\mathbb E[B^{1-z} e^{\eta\mathcal E}]\le 1-(B-1)/8+2 C'\eta\le1
\]
if \(\eta>0\) is sufficiently small depending on \(B\) (and \(\eta\le a_0\)). This gives the required bound on the Rayleigh value and is the induction step for \eqref{cas:e9} when \(k\ge K\).

\medskip\noindent\textbf{The remaining fixed levels.}\enspace
The induction just proved applies when $k\ge K$. For the finitely many
smaller levels, a simpler form of the isolated-path cancellation suffices:
we may allow constants depending on $k$ and bound a Hilbert--Schmidt
norm directly. Proposition~\ref{n:coherent-contact} obtains constants
uniform in $k$ by estimating a squared kernel and using Schur's test.
We give the fixed-$k$ argument here so the base of this induction uses
only harmonic zero sums and the elementary forest count below.
\begin{lemma}[Fixed first-row deficit]\label{cas:fixed-level}
For every fixed integer $k\ge2$ there are constants $C_k,n_k$ such that,
for every dyadic $n=2^d\ge n_k$ and every $\lambda\vdash n$ with
$\lambda_1=n-k$,
\[
 \|K_n(\lambda)\|_{\op}\le C_k(d/n)^{k/2}.
\]
For $k=1$, $K_n(\lambda)=0$ at every admissible size.
\end{lemma}
\begin{proof}
For \(k=1\), the directed sweep sends a single position uniformly over all positions (each new bit is fair conditionally as its layer is reached), so the nontrivial irreducible with \(n-\lambda_1=1\) is annihilated on the one-position module. For fixed \(k\ge2\) and large \(n\), take a function \(f\) in the \(\lambda\) subspace of the ordered distinct \(k\)-tuple module with \(\lambda_1=n-k\). It is orthogonal to functions omitting any slot, by the branching facts (each such entire function space is of types available with \(k-1\) slots). Testing against indicators specifying the other slots, we see its extension by zero to the unrestricted tuple space has zero sum in each slot with the others fixed (also automatically so if the others are not distinct).

From a given distinct \(x\) to a candidate \(y\) in a directed sweep, each proposed path is determined by those two endpoints: updated bits must already equal the corresponding bits of \(y\). For tags \(a,b\), let \(h=h_{ab}\) be the largest differing bit index of \(x_a,x_b\). At layer \(h\) their paths use a shared switch exactly if their \(y\)-bits before \(h\) agree, in which case they require opposite \(y\)-bits at \(h\) for validity. They cannot share at an earlier layer (look at bit \(h\)), or at a later layer on a valid route (their prefixes through \(h\) must then differ). A collision making routing impossible occurs if some such pair instead agrees in \(y\) through \(h\); conversely if this never occurs all paths are on distinct positions at every boundary (compare the unupdated suffix of \(x\) before processing \(h\), or the updated prefix of \(y\) once \(h\) has been processed, for each pair). In that case the needed switch choices are consistent: at a shared switch the two tags distinct before it are sent to distinct outputs of the switch. Each used switch coin is prescribed, with a shared switch reducing the number of prescriptions by one from one-per-tag-per-layer. It follows that the kernel for the directed sweep on these tuples, extended by zero to invalid \(y\), is \(n^{-k}F_{[k]}\), using the functions
\[
F_S=\prod_{a<b:\ a,b\in S}
\left(1-\mathbf1_{\{y_{a,<h_{ab}}=y_{b,<h_{ab}}\}}
                  (-1)^{y_{a,h_{ab}}+y_{b,h_{ab}}}\right).
\]
Here the \(<h\) subscript extracts the indicated prefix, possibly empty. Each factor on a valid route is 2 for a shared switch, 1 otherwise. Against the extended \(f\) summed on the unrestricted tuple space, we can replace \(F_{[k]}\) by \(F'=\sum_{S\subseteq[k]}(-1)^{k-|S|}F_S\) by the zero sums (\(F_S\) uses no \(y\)-slots outside \(S\)). It is bounded by a constant depending on \(k\), and cancels to zero if any vertex is isolated in the graph on tags with potential interaction edges \(y_{a,<h_{ab}}=y_{b,<h_{ab}}\): adding an isolated vertex does not change \(F_S\).

With uniform distinct \(x\) and independent unrestricted uniform \(y\), the probability of no isolates is at most \(C_k(d/n)^{\lceil k/2\rceil}\) for large \(n\). Indeed then the graph contains a spanning forest with no isolates, thus at least \(\lceil k/2\rceil\) edges. For unrestricted independent uniform coordinates also in \(x\), the constraints along any such forest (requiring \(x_a\ne x_b\) along its edges and the indicated prefix matches) cost \((d/n)^{\#\mathrm{edges}}\) by successively integrating leaves. Given the other tags' values at a leaf-removal step, each possible last differing index \(h\in\{1,\ldots,d\}\) relative to the neighbor of the leaf has probability \(2^{h-1}/n\) in \(x\) and the prefix match in \(y\) costs \(2^{-(h-1)}\), totaling \(d/n\). Conditioning \(x\) to be globally distinct costs at most factor 2 for large \(n\) at fixed \(k\), and there are only boundedly many forests at fixed \(k\) in the union bound.

Thus with these tuple measures (distinct for \(x\), unrestricted for \(y\)), \(\mathbb E_{x,y}|F'(x,y)|^2\le C_k(d/n)^{\lceil k/2\rceil}\). The kernel applied to \(f\) at \(x\) is \(\mathbb E_y F'(x,y)\widetilde f(y)\) with \(\widetilde f\) the zero extension. By Cauchy--Schwarz in \(y\), then taking the norm using uniform distinct \(x\), the output norm is at most \(C_k(d/n)^{k/4}\|f\|_2\); here the \(L^2\)-norms use probability measures and the unrestricted norm of \(\widetilde f\) is no larger than \(\|f\|_2\). This kernel on distinct-tuple functions is \(T_n^*\) for the action permuting the tuple basis. Since \(\lambda\) occurs and the averaged permutations preserve its subspace, \(\|T_n^*\|_\lambda=\|T_n\|_\lambda\) gives
\[
\|K_n\|_\lambda\le C_k(d/n)^{k/2}
\]
for each fixed \(k\ge2\) at large \(n\).

\end{proof}

\medskip\noindent\textbf{Choosing the constants and completing the induction.}\enspace
Both the large-level induction and the fixed-level bound have been
proved.  The choices below make them hold with one pair of absolute
constants on every dyadic size.
First fix the integer $K$ from~\eqref{cas:e16}. Since $c/n\le k$ and
$u(2+\log(n/u))$ is increasing for $0<u\le n$,
\[
 \Phi_n(c)\le k(2+\log(n/k)).
\]
For each $1\le k<K$, Lemma~\ref{cas:fixed-level} gives a bound
$o(n^{-k/8})$. There is therefore one absolute $n_0$ such that these
bounds imply~\eqref{cas:e9} whenever $n>n_0$, simultaneously for
all $1<B\le2$ and $0<\eta\le1/8$. Indeed throughout those ranges,
\[
 B^{-1}e^{-\eta\Phi_n(c)}
 \ge\tfrac12 e^{-k/4}k^{k/8}n^{-k/8}.
\]

Only finitely many sizes and shapes remain. On each nontrivial shape
at $n\le n_0$, the norm of $K_n$ is strictly less than one by
Lemma~\ref{lem:finitegap}: equality would give a vector fixed by
every coordinate layer and hence by the hypercube edge transpositions,
which generate $S_n$. Choose $1<B\le2$ close enough to one that
$B^{-1}$ strictly exceeds all these finitely many norms. Then choose
$\eta>0$ small enough to satisfy~\eqref{cas:e9} on this finite set,
$\eta\le\min(1/8,a_0)$, and the large-level induction condition
$2C'\eta\le(B-1)/8$ from~\eqref{cas:e15}.

The direct fixed-level estimates now cover $k<K$ above $n_0$, and the
induction step covers $k\ge K$. This proves~\eqref{cas:e9} on every
dyadic size with one pair of absolute constants $B,\eta$.

\subsection{Hausdorff--Young amplification to the full permutation}

We have separately proved a uniform injection-density bound and an
operator contraction.  To combine them, first convert the density bound
into the $65$th trace moment in a single irreducible.  Then use the
operator contraction for the remaining sweep factors.  The passage to
powers of $T_n$ uses Schatten H\"older, so it does not require $T_n$ to be
normal.

We view probability laws on \(S_n\) also by their densities relative to normalized uniform measure. For a density or function \(f\), use Fourier matrix \(\widehat f(\lambda)=\mathbb E_{g\ {\rm uniform}} f(g) U_\lambda(g)\), with \(U_\lambda\) a unitary irreducible of dimension \(d_\lambda=f^\lambda\). By finite group orthogonality/Plancherel, \(\|f\|_2^2=\sum_\lambda d_\lambda\|\widehat f(\lambda)\|_{\rm HS}^2\) with the (unnormalized) Hilbert--Schmidt norm; this holds with the displayed convention since one is expanding against the complex conjugates of matrix coefficients (scaled by \(\sqrt{d_\lambda}\) for an orthonormal basis). We need the single-irrep Hausdorff--Young bound~\cite[Lemma~5.1(1)]{garcia-cuerva-parcet2004}, using Schatten norms \(S^s\) (the \(\ell^s\)-norms of the singular values, with \(S^\infty\) the operator norm):
\begin{equation}
d_\lambda^{1/s}\|\widehat f(\lambda)\|_{S^s}\le\|f\|_{p},
\qquad 1<p\le2,\quad s=p/(p-1).
\label{cas:e17}
\end{equation}
For completeness, at \(p=1,s=\infty\) the bound without dimension factor is unitarity, and at \(p=s=2\) it follows from Plancherel. To interpolate, pair by trace against a matrix \(A\) of \(S^p\)-norm 1 and take \(\|f\|_p=1\), ignoring the trivial zero case. In \(0\le\operatorname{Re}z\le1\), replace \(f\) by \(f_z\) using its phases times \(|f|^{p(1-z/2)}\) at nonzero entries (keep zeros); similarly in a singular value factorization of \(A\) keep both unitary factors and replace nonzero singular values \(\beta\) by \(\beta^{p(1-z/2)}\) to give \(A_z\). Apply the three-lines theorem to \(d_\lambda^{z/2}\operatorname{Tr}(A_z\widehat f_z(\lambda))\). On the first boundary use \(L^1\to S^\infty\) and the trace norm of \(A_z\), on the second \(L^2\to S^2\) and Hilbert--Schmidt; the norms of the inputs in these bounds are 1 and the analytic family is bounded on the strip. Taking \(z=2(1-1/p)\) and Schatten duality proves \eqref{cas:e17}.

\medskip\noindent\textbf{The injection multiplicity and the $65$th moment.}\enspace
The density estimate controls the Fourier matrix after averaging over
a tuple stabilizer. Averaging these stabilizer projections in turn
recovers $K_n$ at the cost of the ratio $D_\lambda/f^{\bar\lambda}$.
This ratio is at most $\binom nk$, even when the lower-diagram dimension
$f^{\bar\lambda}$ is large. It is this multiplicity comparison that
allows a level-dependent density bound to control the dimension-weighted
trace moment.
Take \(p=1+\delta\) from \eqref{cas:e6}, so \(s=65\). For a fixed nontrivial \(\lambda\), use \(k=n-\lambda_1\). For ordered distinct \(x\), average the density of \(\nu_n\) over right multiplication by its tuple stabilizer \(H_x\). The resulting function is \(g\mapsto R_x(g x)\), with \(L^p\)-norm bounded by \eqref{cas:e6}; its Fourier matrix in \(\lambda\) is \(K_n P_x\), where \(P_x\) is the average of \(U_\lambda\) over \(H_x\). Averaging these projections in \(x\) uniformly gives \((f^{\bar\lambda}/d_\lambda)I\), by conjugation symmetry, Schur's lemma and the fixed-space dimension above. By \eqref{cas:e17} and the triangle inequality in \(S^s\) we have in \(\lambda\)
\[
d_\lambda^{1/s}(f^{\bar\lambda}/d_\lambda)\|K_n\|_{S^s}
=d_\lambda^{1/s}\|\mathbb E_x K_n P_x\|_{S^s}
\le \exp(C k\log(e n/k)).
\]
Taking the \(s\)-th power and using \eqref{cas:e7} gives
\begin{equation}
d_\lambda \operatorname{Tr}_\lambda(K_n^s)
\le \left(\frac{d_\lambda}{f^{\bar\lambda}}\right)^s
       \exp(C s k\log(e n/k))
\le \exp(C_2 k\log(e n/k)).
\label{cas:e18}
\end{equation}
Here \(K_n\) is positive semidefinite and \(C_2\) is absolute (using \(\binom{n}{k}\le(e n/k)^k\)).

\medskip\noindent\textbf{Completion for directed sweeps.}\enspace
After \(l\) directed sweeps with integer \(l>s\), the nontrivial Fourier matrices of the position permutation law/density are \(T_n^l\). Subtracting 1 from the density keeps these and gives zero in the trivial representation. Schatten Hölder for matrix products gives \(\|T_n^s\|_{\rm HS}\le\|T_n\|_{S^{2s}}^s\), so in nontrivial \(\lambda\)
\[
\|T_n^l\|_{\rm HS}^2\le
\|K_n\|_\lambda^{l-s}\operatorname{Tr}_\lambda(K_n^s).
\]
Choose a fixed absolute integer \(l>s\) sufficiently large, with \((l-s)a-C_2\ge 4\) using \(a\) from the consequence of \eqref{cas:e9}. By \eqref{cas:e9}, \eqref{cas:e18} and Plancherel, the squared \(L^2\) distance of the density to 1 is bounded by
\[
\sum_{k=1}^{n-1} 2^k \exp\big(-4 k\log(e n/k)\big)=o(1).
\]
We used that the number of shapes with a given \(k\) is bounded by the number of partitions of \(k\), hence by \(2^k\), and increased \(l\) so the exponent bound applies before counting shapes. The sum tends to zero, for example splitting at \(\sqrt n\) and using \(\log(e n/k)\ge1\) in the large range. Total variation is at most half the \(L^2\) distance by Cauchy--Schwarz. The estimate is independent of the starting permutation, since changing
that permutation translates the law.  As one sweep consists of $d$
physical shuffles, this proves Theorem~\ref{cas:main}.
